\section{Introduction}
\label{sec:intro}

Flood frequency analysis (FFA) rests on the assumption of stationarity: that the sequence
of annual maximum discharges is drawn from a single, time-invariant probability
distribution. When a catchment is affected by a recurring process that inflates extreme
flows in a subset of years, this assumption is violated, and fitted design quantiles can
be misleading.

The Simme river (canton Bern, Switzerland) provides a compelling case study. The Plaine
Morte glacier hosts a supraglacial lake (Lac des Faverges) whose drainage has produced
documented Glacial Lake Outburst Floods (GLOFs) in 2011, 2012, 2013, 2014, 2015, 2017,
and 2018. Each event releases a sub-daily flood pulse that consistently ranks among the
largest annual flows recorded at the gauge. Since GLOFs are mechanistically distinct from
the rainfall- and snowmelt-driven floods that dominate the pre-2011 record, their inclusion
in a stationary FFA is conceptually problematic.

Following the largest recorded GLOF in 2018, which caused approximately CHF 2.5 million in
damage \citep{GeoTest2019}, a supraglacial bypass channel was excavated in summer 2019 to
divert the lake's outflow in a controlled manner through existing englacial drainage
systems. A second parallel channel was added in June 2022 to improve redundancy
\citep{GeoPraevent2022}. These engineering interventions effectively ended uncontrolled
outburst flooding after 2018, rendering the GLOF period 2011--2018 a closed historical
episode. Understanding the statistical imprint of those seven events on the long-term flood
frequency record therefore remains relevant for retrospective hazard assessment and for
evaluating whether the pre-intervention design quantiles were systematically biased.

This study addresses two research questions:
\begin{enumerate}
  \item Do recurring GLOFs systematically shift the GEV distribution fitted to the Simme
        AMS, and if so, by how much?
  \item Are the flood magnitudes observed since 2011 statistically consistent with the
        catchment's natural (pre-GLOF) hydrology once the GLOF signal is removed?
\end{enumerate}
